% Part: set-theory
% Chapter: story
% Section: russells-paradox-again

\documentclass[../../../include/open-logic-section]{subfiles}

\begin{document}

\olfileid{sth}{story}{rus}
\olsection{রাসেলের কূটাভাস (আবার)}

\olref[sfr][][]{part}-এ আমরা একটি অনানুষ্ঠানিক সেটতত্ত্ব নিয়ে কাজ
করেছি। কিন্তু \emph{অত্যন্ত} সরল একটি ধারণা অনুযায়ী, সেট বলতে
কেবল বিধেয়ের বিস্তারকেই বোঝায়। এই সরল ধারণাটি মেনে নিলে
নিচের নীতি গ্রহণ করতে হয়:

\begin{defish}
  \emph{অনানুষ্ঠানিক ধর্মনির্দেশে সেট-গঠন।} যেকোনো সূত্র~$\phi$-এর জন্য
  $\Setabs{x}{\phi(x)}$-এর অস্তিত্ব আছে।
\end{defish}

নীতিটি আকর্ষণীয় হলেও প্রমাণ করা যায় যে এটি অসঙ্গত।
\olref[sfr][set][rus]{sec}-এ আমরা তা দেখেছি। কিন্তু ফলটি
এত গুরুত্বপূর্ণ এবং এত সরল যে, হুবহু আবার বলা দরকার।

\begin{thm}[রাসেলের কূটাভাস]
$R = \Setabs{x}{x \notin x}$ এমন কোনো সেট নেই।
\end{thm}

\begin{proof}
যদি $R = \Setabs{x}{x \notin x}$ থাকে, তবে $R \in R$
যদি এবং কেবল যদি $R \notin R$; এটি স্ববিরোধ।
\end{proof}

রাসেল ১৯০১ সালের জুনে এই ফলটি আবিষ্কার করেন। (তবে আমরা যেভাবে
কূটাভাসটি বললাম, তিনি ঠিক সেভাবে বলেননি; তাঁর বিবেচনায় ছিল
\emph{Grundgesetze}-তে বর্ণিত ফ্রেগের সেটতত্ত্ব।
\olref[blv]{sec}-এ আমরা এ প্রসঙ্গে ফিরব।) ফ্রেগের ব্যবস্থায়
অসঙ্গতিটি ব্যাখ্যা করে রাসেল তাঁকে ১৯০২ সালের ১৬ জুন চিঠি লেখেন।
এই চিঠিপত্র ও কিছু পটভূমির জন্য
\citet[pp.~124--8]{Heijenoort1967} দেখুন।

জোর দিয়ে বলা দরকার যে, এই দুই-সারির প্রমাণটি \emph{বিশুদ্ধ
যুক্তিবিদ্যা}-র ফল। অবশ্য $\Setabs{x}{x \notin x}$ প্রতীকে
আমরা প্রচ্ছন্নভাবে সদস্যভিত্তিক সমতার একটি (যুক্তিবহির্ভূত?)\
স্বতঃসিদ্ধ ব্যবহার করেছি: $\Setabs{x}{\phi(x)}$ বলতে, যদি
এমন সেট থাকে, $\phi$ ধর্মবিশিষ্ট বস্তুগুলির \emph{একমাত্র}
সেটটিকে বোঝায়; একমাত্র হওয়া সদস্যভিত্তিক সমতা থেকে আসে।
কিন্তু ফলটি এভাবে বললে সেই নীতির ইঙ্গিতটুকুও এড়ানো যায়:
\emph{যে সেটগুলির নিজের সদস্য হওয়ার ধর্ম নেই, তাদের ঠিক সকলকে
সদস্য হিসেবে ধারণ করে এমন কোনো সেট নেই}। এর সঙ্গে সেটের
বিশেষ কোনো যোগও নেই। রাসেল নিজেই লক্ষ করেছিলেন, ঠিক একই
যুক্তিতে বলা যায়: \emph{যারা নিজের দাড়ি নিজেরা কামায় না,
ঠিক তাদের সকলের দাড়ি কামায় এমন কোনো মানুষ নেই}। অথবা:
\emph{যে পাগ-কুকুরগুলি নিজেদের শোঁকে না, ঠিক তাদের সকলকে শোঁকে
এমন কোনো পাগ-কুকুর নেই}। এভাবেই আরও উদাহরণ দেওয়া যায়।
ছক হিসেবে ফলটির রূপ কেবল:
\[
\lnot \exists x \forall z(Rzx \liff \lnot R zz).
\]
এটি প্রথম-ক্রমের যুক্তিবিদ্যারই একটি উপপাদ্য (বা উপপাদ্য-ছক)।
ফলে শুধু সেটতত্ত্বে সামান্য রদবদল করে রাসেলের কূটাভাস এড়ানো
যাবে না; সেটতত্ত্বে \emph{পৌঁছনোর আগেই} কূটাভাসটি দেখা দেয়।
ধ্রুপদী প্রথম-ক্রমের যুক্তিবিদ্যা ব্যবহার করলে
$R = \Setabs{x}{x\notin x}$ এমন কোনো সেট নেই—এ কথা
আমাদের \emph{মেনে নিতেই হবে}।

ফল দাঁড়ায় এই: অসঙ্গতি \emph{এড়িয়েও} যদি অনানুষ্ঠানিক ধর্মনির্দেশে
সেট-গঠন মেনে নিতে চান, শুধু \emph{সেটতত্ত্ব} বদলালে চলবে না;
\emph{যুক্তিবিদ্যা}-ই আমূল বদলাতে হবে।

অবশ্য অধ্রুপদী যুক্তিবিদ্যার সঙ্গে সেটতত্ত্বও প্রস্তাব করা হয়েছে।
কিন্তু সেগুলি অন্তত প্রচলিত পদ্ধতি নয়। রাসেলের কূটাভাস নিয়ে
প্রচলিত পদ্ধতি হলো একে সরাসরি অনস্তিত্বের প্রমাণ হিসেবে গ্রহণ
করা এবং তারপর এই ফলের সঙ্গে কীভাবে কাজ করা যায়, তা শেখা।
আমরাও সেই পদ্ধতিই অনুসরণ করব।

\end{document}
